\documentclass[11pt,a4paper]{article}
\usepackage[margin=2.3cm]{geometry}
\usepackage{amsmath,amssymb,mathtools,bm}
\usepackage{booktabs}
\usepackage[T1]{fontenc}
\usepackage{lmodern}
\usepackage{microtype}
\usepackage{enumitem}
\usepackage{xcolor}
\usepackage[colorlinks=true,linkcolor=blue!50!black,citecolor=blue!50!black,urlcolor=blue!50!black]{hyperref}

\newcommand{\bphi}{\bm{\phi}}
\newcommand{\bpsi}{\bm{\psi}}
\newcommand{\bc}{\bm{c}}
\newcommand{\cbar}{\bar{\bm{c}}}
\newcommand{\bg}{\bm{g}}
\newcommand{\tran}{^{\mathsf T}}
\newcommand{\E}{\mathbb{E}}
\newcommand{\Prob}{\mathbb{P}}
\newcommand{\diag}{\operatorname{diag}}
\newcommand{\Tfit}{T_{\mathrm{fit}}}
\newcommand{\Tval}{T_{\mathrm{val}}}
\newcommand{\vcal}{v^{\mathrm{cal}}}
\newcommand{\chiinv}{\chi^{-1}_3}

\title{Calibrated per-atom force uncertainty for linear ACE:\\
\large the full mathematical pipeline (current, and the proposed conformal scale)}
\author{ace-jax \texttt{feat/conformal-uq} --- for review}
\date{1 October 2026}

\begin{document}
\maketitle

\noindent Companion to \texttt{docs/specs/2026-09-30-conformal-force-sigma-design.md} (interfaces)
and \texttt{2026-09-28-tempered-ard-uq-design.md} (ARD and sandwich, merged in \#18).
Notation: atoms $i$, Cartesian components $\alpha\in\{x,y,z\}$, training configurations $c\in T$.

\paragraph{Summary: the conformal step does not replace the sandwich.}
The served force uncertainty is a product of two factors.
\begin{itemize}[nosep]
  \item The \emph{shape} $v(x)$ is the unscaled per-atom variance (the sandwich by default). It
    decides how $\sigma$ varies from atom to atom, so it carries all the ranking information. This
    proposal leaves it \textbf{unchanged}.
  \item The \emph{scale} is currently one scalar $\lambda$ (or $\kappa$), fitted so that the rms
    normalised error is~1. The proposal replaces it with one conformal factor $\lambda_g$ per group
    $g$ of a geometric feature, obtained from score \emph{quantiles}. The score itself is normalised
    by the sandwich variance.
\end{itemize}
Within a group, atoms are ordered by $\sigma$ exactly as the sandwich orders them. Across groups the
ordering changes only through the ratios $\lambda_g/\lambda_{g'}$.

\tableofcontents

%======================================================================
\section{Model and data}

Each atom $i$ with environment $x_i$ has a site descriptor $\bphi(x_i)\in\mathbb R^L$: the ACE product
basis in species-major blocks, plus pair blocks. The model is linear in the coefficients $\bc$:
\begin{equation}
  E(\mathbf R) = \sum_i \bphi(x_i)\cdot\bc + \sum_i E^0_{Z_i}.
\end{equation}
A training configuration contributes energy, force and virial observations, each linear in $\bc$:
\begin{equation}
  \bphi_E=\sum_i\bphi(x_i),\qquad
  \bphi_{F,i\alpha} = -\frac{\partial\bphi_E}{\partial r_{i\alpha}},\qquad
  \bphi_{V,\alpha\beta} = -\sum_i r_{i\alpha}\frac{\partial\bphi_E}{\partial r_{i\beta}}.
\end{equation}
Each row has a structural weight $w$ (per-atom energy normalisation and per-type weights) and a
noise scale $\sigma_q$ for its quantity $q\in\{E,F,V\}$. The whitened rows and targets are
\begin{equation}
  \bpsi = \bphi\, w/\sigma_q,\qquad \tilde y = y\, w/\sigma_q,
\end{equation}
where $y$ is the label minus the $E^0$ baseline.

%======================================================================
\section{Prior and posterior}

The prior is $\bc\sim\mathcal N(0,\Lambda^{-1})$ with
\begin{equation}
  \Lambda=\diag\big(\Gamma_j^2\,e^{a_{k(j)}}\big).
\end{equation}
Here $\Gamma_j$ is the smoothness prior of column $j$, $k(j)\in\{2,3,4\}$ is the column's body
order, and $a_k$ is the ARD log-precision multiplier of body order $k$. Pair columns and
correlation-order-1 columns both count as 2-body. The posterior is Gaussian:
\begin{equation}
  A=\sum_{\text{rows}}\bpsi\bpsi\tran+\Lambda,\qquad
  \cbar = A^{-1}\sum_{\text{rows}}\bpsi\,\tilde y,\qquad
  \operatorname{Cov}(\bc)=A^{-1}.
\end{equation}
Numerically we work in the prior-scaled system $S=D^{-1}AD^{-1}=LL\tran$, with $D=\diag\Gamma$.
Each $a_k$ is floored, $a_k\ge a_{\mathrm{floor}}$, so that $\operatorname{cond}(S)\le10^{14}$.

%======================================================================
\section{Hyperparameters: type-II maximum likelihood (unchanged)}

The hyperparameters are $h=(\log\sigma_E,\log\sigma_F,\log\sigma_V,a_2,a_3,a_4)$. The evidence uses
the per-quantity statistics $G_q=\sum\bpsi\bpsi\tran$, $\bm b_q$, $\bm y\tran\bm y_q$ and $n_q$, and
up to constants is
\begin{equation}
  \log p(\mathcal D\mid h) = -\tfrac12\sum_q\frac{\bm y\tran\bm y_q}{\sigma_q^2}
  +\tfrac12\bm b\tran A^{-1}\bm b-\tfrac12\log|A|+\tfrac12\log|\Lambda|-\sum_q n_q\log\sigma_q .
\end{equation}
The \emph{joint} mode maximises it over all six components of $h$, with L-BFGS-B on the
gradient-normalised objective. The \emph{sequential} mode fixes $\sigma_q$ at the linear MAP and
fits only the $a_k$.

%======================================================================
\section{The hold-out protocol (unchanged; it feeds every scale)}
\label{sec:holdout}
\begin{enumerate}[nosep]
  \item Split the training configurations, by configuration, into $\Tfit$ (fraction $1-f$) and
    $\Tval$ (fraction $f=0.2$).
  \item Fit $h$ on $\Tfit$ to obtain the posterior $P_{\mathrm{fit}}$. For each held-out atom
    $i\in\Tval$, record the error $\bm e_i=\bm F_i-\hat{\bm F}_i(\cbar_{\mathrm{fit}})$. These errors
    are \emph{honest}, because $P_{\mathrm{fit}}$ never saw $\Tval$.
  \item Refit on all of $T$, starting from $h_{\mathrm{fit}}$. This gives the \textbf{served}
    posterior $P$, with mean $\cbar$ and factor $L$.
\end{enumerate}

%======================================================================
\section{The shape \texorpdfstring{$v(x)$}{v(x)}: unscaled per-atom variance (unchanged)}
\label{sec:shape}

A target atom has force rows $\bphi_\alpha(x)$, $\alpha=x,y,z$ (raw, unweighted). There are two
variants of the shape.

\paragraph{Epistemic (\texttt{-{}-ard-variance kappa}).}
\begin{equation}
  s^2(x)=\sum_\alpha\bphi_\alpha A^{-1}\bphi_\alpha\tran=\sum_\alpha\big\|L^{-1}D^{-1}\bphi_\alpha\tran\big\|^2.
\end{equation}

\paragraph{Configuration-clustered sandwich (\texttt{-{}-ard-variance sandwich}, default).}
The residuals at the served mean are $\rho_r=\tilde y_r-\bpsi_r\cdot\cbar$. Each training
configuration's score sums its rows, over all of $c$'s energy, force and virial rows:
\begin{equation}
  \bg_c=\sum_{r\in c}\bpsi_r\rho_r .
\end{equation}
The ``meat'' is $M=\sum_c\bg_c\bg_c\tran$, of rank at most $|T|$, and
\begin{equation}
  m^2(x)=\sum_\alpha\bphi_\alpha A^{-1}MA^{-1}\bphi_\alpha\tran
  =\sum_\alpha\big\|Q\tran D^{-1}\bphi_\alpha\tran\big\|^2,
  \qquad Q=S^{-1}D^{-1}\big[\bg_1\cdots\bg_{|T|}\big]\in\mathbb R^{L\times|T|}.
  \label{eq:sandwich}
\end{equation}
This is the Huber--White misspecification-robust covariance of $\hat{\bc}$, and also M\"uller's
(2013) sandwich posterior; $\kappa^2A^{-1}$ is its homoscedastic special case ($M\propto A$). The
term $\bphi_\alpha A^{-1}\bg_c$ is the one-step change in the prediction at $x$ when configuration
$c$ is deleted. So $m^2(x)$ is the \emph{infinitesimal-jackknife} variance of the prediction over
training configurations. This is what ranks local errors: on bench365 the Spearman correlation
$\rho$ is $0.30$--$0.37$ for $m^2$, against $0.19$--$0.26$ for $s^2$.

\paragraph{Leave-own-cluster-out (calibration only).}
A held-out atom $i$ belongs to training configuration $c(i)$, and a genuinely new configuration has
no ``own'' term. For calibration we therefore use
\begin{equation}
  m^2_{-\mathrm{own}}(x_i)=\sum_\alpha\sum_{c\neq c(i)}\big(\bphi_\alpha A^{-1}\bg_c\big)^2 .
\end{equation}
The \emph{served} variances use all $c$. Below, $v(x)$ denotes the chosen shape ($m^2$ or $s^2$), and
$\vcal_i$ its calibration version ($m^2_{-\mathrm{own}}$, or $s^2$) at the held-out atoms.

%======================================================================
\section{The scale: current (Gaussian) and proposed (conformal)}
\label{sec:scale}

For a labelled calibration atom, the normalised score is
\begin{equation}
  s_i=\frac{|\bm e_i|}{\sqrt{\vcal_i/3}} .
\end{equation}
If $\sqrt{v/3}$ were an exact isotropic Gaussian per-component scale, $s$ would follow $\chi_3$, the
chi distribution with 3 degrees of freedom.

\subsection{Current: Gaussian rms matching (one scalar)}
\begin{equation}
  \lambda^2=\frac13\,\overline{s^2}_{\,i\in\Tval}\qquad(\text{the Gaussian NLL optimum; $\kappa$ is the same formula with $s^2$}).
\end{equation}
The served value is $\sigma(x)=\lambda\sqrt{v(x)}$, read as $\Delta\bm F\sim\mathcal N(0,\sigma^2/3\,I_3)$.
The only claim made is rms-$z=1$ on $\Tval$; any coverage statement rests on the Gaussian
assumption.

\subsection{Proposed: Mondrian split conformal (one factor per group)}

The groups $g(x)\in\{1,\dots,G\}$ are computed from geometry alone (Section~\ref{sec:groups}). For a
target coverage $1-\alpha$, and each group $g$ with $n_g$ calibration atoms, take
\begin{equation}
  q_g=\text{the }\frac{\lceil(n_g+1)(1-\alpha)\rceil}{n_g}\text{ empirical quantile of }\{s_i: g(x_i)=g\}.
\end{equation}
The served quantities are
\begin{align}
  \texttt{forces\_q}(x)  &= q_{g(x)}\sqrt{v(x)/3}, \label{eq:fq}\\
  \texttt{forces\_std}(x)&= \lambda_{g(x)}\sqrt{v(x)},\qquad \lambda_g\coloneqq q_g/\chiinv(1-\alpha). \label{eq:fs}
\end{align}
Here \eqref{eq:fq} is the guaranteed radius, and \eqref{eq:fs} is the $\sigma$ whose Gaussian
$(1-\alpha)$ ball is that radius.

\paragraph{Guarantee (split conformal, per group).}
If a target atom is exchangeable with group $g$'s calibration atoms, then
\begin{equation}
  1-\alpha\;\le\;\Prob\big(|\Delta\bm F|\le\texttt{forces\_q}\big)\;\le\;1-\alpha+\frac{1}{n_g+1}.
\end{equation}
No Gaussianity is assumed: within a group, $s$ may have any distribution. Atoms within a
configuration are correlated, so the exact statement holds at the level of configurations, and the
atom-level version is the standard approximation. We report $n_{\mathrm{cfg}}$ per group and merge
groups below $n_{\min}$.

\paragraph{Relation to the Gaussian scale.}
The Gaussian scale is the case $G=1$, with the rms in place of the $(1-\alpha)$ quantile. In
distribution the two scales agree to about 1\%: on bench365, $q=2.48$ against
$\lambda\,\chiinv(0.9)=2.50$. Where they differ, the conformal scale is the one with a stated coverage.

\paragraph{What is invariant.}
The mean $\cbar$ and the energy and virial predictions; the shape $v(x)$; and the ordering of atoms by
$\sigma$ \emph{within a group}, since the transformation there is a single positive factor.
What changes is the relative $\sigma$ \emph{between} groups, by $\lambda_g/\lambda_{g'}$. This is how
crack-tip groups receive more $\sigma$ without inflating the bulk.

%======================================================================
\section{The group feature (Mondrian)}
\label{sec:groups}

Two constants are fixed at fit time: $r_1$, the first minimum of the training radial distribution
function, and $z^\star$, the modal training coordination. For each atom,
\begin{equation}
  z_i=\#\{j:r_{ij}<r_1\},\qquad
  d_i=\frac{\operatorname{std}\{r_{ij}:r_{ij}<r_1\}}{\operatorname{mean}\{r_{ij}:r_{ij}<r_1\}}
  \quad(\text{NaN if }z_i<2).
\end{equation}
The band edges are quantiles of $d$ over the calibration pool: $p_{50}$, $p_{90}$ and $p_{99}$ by
default, giving four bands. The group is
\begin{equation}
  g=\operatorname{band}(d)\times[\,z=z^\star\,],\qquad G=8 .
\end{equation}
These features were chosen from the big-cell results. There, miscoverage tracked distortion ($\sigma$
grows $\times1.7$ from low to high distortion, while the error grows $\times2.0$) and coordination:
bulk-like atoms were calibrated, surface-like atoms conservative, and atoms with a strained first
shell overconfident. In distribution, coverage is flat across these groups, so the grouping costs
calibration samples only where nothing is wrong.

%======================================================================
\section{Calibration sets}

\paragraph{Default (inside the fit).}
$\Tval$, with its honest errors (Section~\ref{sec:holdout}) and $\vcal$ (Section~\ref{sec:shape}).

\paragraph{\texttt{aj calibrate} (target regime).}
A set $U$ of labelled configurations, for example a few cells from the production regime labelled
with MACE or DFT. The errors are those of the \emph{served} model, $\bm e_u=\bm F_u-\hat{\bm
F}_u(\cbar)$; since $U$ is disjoint from training, these are honest too. The variance $v_u$ is the
served $v$, with no own term because $U\notin T$. The scores $s_u$ are appended to $\Tval$'s within
each group (or replace them, with \texttt{-{}-replace}), and the $q_g$ are recomputed.
Exchangeability is then with $\Tval\cup U$ within each group, so target atoms that resemble $U$ are
covered because $U$ is in the pool.

\paragraph{Evidence.}
On bench365, leaving out one realisation at a time over 10 crack realisations (30 cells): the whole
crack cell reaches coverage $0.899$, the crack tip $0.885$ (fold range $0.864$--$0.900$), and the
dislocations $0.92$--$0.93$.

%======================================================================
\section{The support diagnostic (covariate-shift weighted conformal; not served)}

Our labels are a deterministic function of the structure, so $P(y\mid x)$ is fixed and going from
training to a target is a pure \emph{covariate shift}. For a target structure, given only its
unlabelled atoms:
\begin{enumerate}[nosep]
  \item \textbf{Density ratio.} For each species, fit an L2 logistic classifier on whitened 2-body
    descriptors, with target $=1$ and calibration $=0$, class-balanced. This gives
    $w(x)\propto P(t\mid x)/P(c\mid x)$.
  \item \textbf{Weighted quantile} (Tibshirani, Barber, Cand\`es \& Ramdas, 2019). Let
    $W=\sum_j w(x_j)+w(x)$. For each target atom $x$,
    \begin{equation}
      q_w(x)=\text{the }(1-\alpha)\text{ quantile of }
      \sum_{i\in\mathrm{cal}}\frac{w(x_i)}{W}\,\delta_{s_i}+\frac{w(x)}{W}\,\delta_{+\infty}.
    \end{equation}
  \item \textbf{Report} \texttt{support\_q}$=q_w(x)$ and \texttt{support\_ok}$=[q_w<\infty]$ per
    atom, and $n_{\mathrm{eff}}=(\sum w)^2/\sum w^2$ per species.
\end{enumerate}
A value \texttt{support\_ok}$=$False means that, under the estimated shift, the calibration data
cannot certify coverage for that atom; these are the atoms to label next (active learning). On
bench365 with no target labels, the check flags $10.7\%$ of crack-tip atoms against $0.4$--$0.5\%$
for the dislocations. The diagnostic is \emph{not} used to scale $\sigma$. With target labels in the
pool it adds nothing over the Mondrian scale, and without them its finite quantiles rest on
$n_{\mathrm{eff}}\sim10$--$10^3$.

%======================================================================
\section{End-to-end summary}

\begin{center}
\fbox{\begin{minipage}{0.93\linewidth}
\begin{description}[leftmargin=2.2cm,style=nextline]
\item[fit] Data $\to$ whitened rows $\bpsi$, $\tilde y$ $\to$ type-II ML $h$ (on $\Tfit$; refit on $T$)
  $\to$ served posterior $(\cbar,L)$ $\to$ configuration scores $\bg_c$ at $\cbar$ $\to$
  $Q=S^{-1}D^{-1}G$ (the shape) $\to$ held-out scores $s_i=|\bm e_i|/\sqrt{\vcal_i/3}$ with groups
  $g(x_i)$ $\to$ per-group $q_g$ (conformal) or a scalar $\lambda$ (Gaussian).
\item[calibrate] A labelled target set $U$ $\to$ append $s_u$ within each group $\to$ recompute $q_g$.
\item[serve] $v(x)$ (the shape) $\times\,\lambda_{g(x)}=q_{g(x)}/\chiinv(1-\alpha)$ $\to$
  \texttt{forces\_std}, \texttt{forces\_q}.
\item[diagnose] An unlabelled target $\to$ $w(x)$ $\to$ weighted quantile $\to$
  \texttt{support\_ok}, \texttt{support\_q}, $n_{\mathrm{eff}}$.
\end{description}
\end{minipage}}
\end{center}

%======================================================================
\section{Assumptions and known limitations}
\begin{enumerate}[nosep]
  \item \textbf{Exchangeability within a group, at configuration level.} The per-atom guarantee is
    approximate; the groups and $n_{\mathrm{cfg}}$ are reported.
  \item \textbf{The groups were chosen from benchmark evidence.} A shift the groups do not resolve
    (chemical short-range order, a new phase) is covered only marginally. The support diagnostic is
    the guard: it flags unsupported atoms without needing labels.
  \item \textbf{A residual gap at the crack tip.} Even with labelled crack cells, coverage at the tip
    itself is about 1.5 points below target. No feature tried (distortion, coordination, 2-body
    density ratio, 2-body $k$NN) resolves the tip completely.
  \item \textbf{Forces only.} The energy and virial variances remain the unscaled posterior
    variances, unchanged.
  \item \textbf{No aleatoric term.} The labels are deterministic, so all error is approximation
    error. The scale ($\lambda$ or $q_g$) absorbs its size; the shape $v(x)$ decides where it is
    larger.
\end{enumerate}

\begin{thebibliography}{9}\small
\bibitem{huber} P.~J.~Huber, \emph{The behavior of maximum likelihood estimates under nonstandard
  conditions}, Proc.\ 5th Berkeley Symp.\ (1967); H.~White, Econometrica 50, 1 (1982).
\bibitem{muller} U.~K.~M\"uller, \emph{Risk of Bayesian inference in misspecified models, and the
  sandwich covariance matrix}, Econometrica 81, 1805 (2013).
\bibitem{tibshirani} R.~J.~Tibshirani, R.~F.~Barber, E.~J.~Cand\`es, A.~Ramdas, \emph{Conformal
  prediction under covariate shift}, NeurIPS (2019), arXiv:1904.06019.
\bibitem{vovk} V.~Vovk, A.~Gammerman, G.~Shafer, \emph{Algorithmic Learning in a Random World},
  Springer (2005) --- Mondrian conformal prediction.
\end{thebibliography}

\end{document}
